\documentclass[11pt]{article}
\usepackage[margin=1in]{geometry}
\usepackage{amsmath,amssymb,booktabs,microtype,xcolor}
\usepackage[colorlinks=true,linkcolor=blue!50!black]{hyperref}
\usepackage{caption}
\captionsetup{font=small}
\newcommand{\E}{\mathbb{E}}
\title{\textbf{Which statistics give each generator away?}\\[2pt]
\large Energy distance one statistic at a time, on the fixed engine (R36i)}
\author{ProductGPT technical note (R36i)}
\date{31 October 2026}

\begin{document}
\maketitle

\begin{abstract}
\noindent The joint energy distance says how far a generator's holdout
histories land from the real ones. This note breaks it down by statistic: one
energy distance, with its own real-vs-real band, for each of the 27 statistics,
for every arm of R36i. Free running fails almost everywhere, mainly by
overspending. Scheduled sampling halves the damage but keeps the same pattern.
The Markov chain gets the levels right and misses memory beyond one step.
Teacher forcing misses run lengths and switching. Arm C fails on one statistic
only, the satiation signature at one copy, which locates the lottery's known
residual and qualifies that statistic for every arm that uses the lottery.
\end{abstract}

\section{Data and definition}

\paragraph{Source.} R36i, job 42879 (20 October 2026), on the fixed rollout
engine: 300 out-of-sample customers, holdout campaigns 28--30, eight replicates
per customer, three runs per arm (training seeds for model arms, draw seeds
otherwise). The arms are those of the generative-failure figure: free running
of the b33 model, scheduled sampling at $p = 0.15$ and $0.35$ (b35), the
first-order Markov chain, independent draws, and the two diagnostic controls,
arm C and teacher forcing. The controls read the realised holdout, so they
locate failures but cannot simulate a policy. The saved per-customer statistics
reproduce all 21 joint energies in the ledger to within $10^{-15}$.

\paragraph{Definition.} For statistic $k$, with real and generated distribution
functions $F_k$ and $G_k$,
\begin{equation}
\mathcal{E}_k \;=\; 2\,\E|X_k - Y_k| - \E|X_k - X_k'| - \E|Y_k - Y_k'|
\;=\; 2\int_{-\infty}^{\infty}\big(F_k(x) - G_k(x)\big)^2\,dx .
\end{equation}
Each statistic is standardised by the real mean and standard deviation, and
undefined values are set to the real mean, exactly as for the joint distance.
The estimator, the matched half-samples (150 against 150) and the 120
split-half draws are also the same, so each statistic's band is built on the
same customer splits as the joint band (0.078).

\paragraph{How to read the numbers.} Four cautions.
\begin{enumerate}
\item \textbf{These tables cannot see dependence.} A one-dimensional distance
  is blind to how statistics move together. A generator could pass all 27 and
  still fail jointly.
\item \textbf{Sparse statistics have tiny bands.} The satiation signature at
  $k = 2$ and $k = 3{+}$ is undefined for 89\% and 97\% of real customers (few
  hold two or more copies of a featured product while it is on offer). Their
  bands are therefore 0.003 and 0.001, and their ratios mostly reflect how often
  the statistic is defined, which the ``undefined'' column shows.
\item \textbf{27 statistics, one band each.} An exact generator would still land
  outside on about one statistic in twenty. Read patterns, not single
  exceedances.
\item \textbf{The lottery's own error.} Arm C, which uses real decisions and the
  simulated lottery, fails on the satiation signature at one copy (Section~2).
  Every arm that draws acquisitions from the lottery inherits that error.
\end{enumerate}

\section{The overview}

Table~\ref{tab:overview} expresses each arm's mean energy as a multiple of that
statistic's band. Values above 1 are outside on average; grey values are inside
on all three runs.

\begin{table}[h]
\centering
\small
\caption{Mean energy over three runs, in multiples of each statistic's
real-vs-real band. Grey: inside the band on every run. Bottom row: statistics
outside the band on every run.}
\label{tab:overview}
% Generated by scripts/r36i_per_statistic.py -- do not edit by hand.
\begin{tabular}{@{}lrrrrr|rr@{}}
\toprule
 & \multicolumn{5}{c|}{generators} & \multicolumn{2}{c}{controls}\\
statistic & free run. & SS 0.15 & SS 0.35 & Markov & i.i.d. & arm C & teacher f.\\
\midrule
Spend & 27.8 & 12.8 & 8.5 & \textcolor{gray}{0.2} & \textcolor{gray}{0.4} & \textcolor{gray}{-0.2} & 2.6\\
Purchase incidence & 12.4 & 6.0 & 4.2 & 1.9 & 2.7 & \textcolor{gray}{-0.1} & 0.9\\
Share of Buy1, regular banner & 5.2 & 3.0 & 3.6 & 3.4 & 4.7 & \textcolor{gray}{-0.2} & \textcolor{gray}{0.0}\\
Share of Buy10, regular banner & 19.8 & 20.1 & 7.7 & 9.5 & 12.3 & \textcolor{gray}{-0.1} & 8.6\\
Share of Buy1, figure A & 3.5 & 3.3 & 2.4 & 2.1 & 4.0 & \textcolor{gray}{-0.1} & \textcolor{gray}{-0.0}\\
Share of Buy10, figure A & 1.3 & \textcolor{gray}{0.3} & 0.6 & 1.1 & 3.0 & \textcolor{gray}{-0.2} & 0.9\\
Share of Buy1, figure B & 4.6 & 5.3 & 7.8 & 12.1 & 19.9 & \textcolor{gray}{-0.2} & \textcolor{gray}{0.4}\\
Share of Buy10, figure B & 53.1 & 36.0 & 18.4 & 4.0 & 9.2 & \textcolor{gray}{-0.1} & 9.4\\
Share of Buy1, weapon banner & 21.3 & 18.0 & 17.4 & 2.6 & 7.2 & \textcolor{gray}{-0.1} & 1.0\\
Share of Buy10, weapon banner & 179.1 & 77.1 & 71.9 & 1.0 & 7.1 & \textcolor{gray}{-0.1} & 4.6\\
Share of NotBuy & 16.7 & 8.0 & 5.6 & 2.4 & 3.6 & \textcolor{gray}{-0.1} & 1.1\\
Longest NotBuy run & 20.6 & 11.7 & 9.6 & 15.8 & 27.6 & \textcolor{gray}{-0.1} & 9.2\\
Mean buying run & 28.6 & 6.9 & 6.2 & \textcolor{gray}{0.9} & 9.9 & \textcolor{gray}{-0.1} & 0.8\\
Longest buying run & 14.8 & 4.6 & 4.1 & 1.1 & 6.7 & \textcolor{gray}{-0.1} & \textcolor{gray}{-0.1}\\
Banner switch rate & 1.2 & \textcolor{gray}{-0.0} & \textcolor{gray}{0.4} & 7.0 & 78.9 & \textcolor{gray}{-0.2} & 16.3\\
Spend concentration & 1.0 & 1.1 & \textcolor{gray}{0.5} & 4.5 & 7.9 & \textcolor{gray}{-0.1} & 2.0\\
Autocorrelation, lag 1 & 23.2 & 12.3 & 8.4 & 11.7 & 101.1 & \textcolor{gray}{-0.1} & 17.0\\
Autocorrelation, lag 2 & 19.2 & 11.3 & 7.6 & 31.8 & 77.2 & \textcolor{gray}{-0.1} & 8.8\\
Autocorrelation, lag 3 & 17.0 & 10.0 & 6.9 & 41.9 & 63.2 & \textcolor{gray}{-0.2} & 5.9\\
Autocorrelation, lag 4 & 14.5 & 8.6 & 5.9 & 39.6 & 48.1 & \textcolor{gray}{-0.2} & 4.7\\
Autocorrelation, lag 5 & 9.9 & 5.9 & 4.2 & 29.5 & 32.3 & \textcolor{gray}{-0.1} & 2.2\\
Escalation & 7.6 & 1.4 & 1.0 & \textcolor{gray}{0.3} & 43.9 & \textcolor{gray}{-0.2} & 6.0\\
Persistence & 22.3 & 9.5 & 8.6 & \textcolor{gray}{0.8} & 45.4 & \textcolor{gray}{-0.2} & 2.7\\
Satiation signature, $k = 0$ & 18.6 & 9.1 & 7.9 & 2.7 & 5.4 & \textcolor{gray}{0.7} & 1.9\\
Satiation signature, $k = 1$ & 57.6 & 36.8 & 33.8 & 7.4 & 6.9 & 7.7 & 1.9\\
Satiation signature, $k = 2$ & 26.1 & 8.6 & 7.5 & 1.6 & 1.7 & \textcolor{gray}{0.9} & \textcolor{gray}{-0.1}\\
Satiation signature, $k = 3+$ & 22.4 & 6.8 & 5.1 & \textcolor{gray}{0.4} & \textcolor{gray}{0.4} & \textcolor{gray}{0.3} & \textcolor{gray}{-0.1}\\
\midrule
outside on every run & 24 & 24 & 15 & 20 & 25 & 1 & 16\\
\bottomrule
\end{tabular}

\end{table}

\paragraph{What it shows.}
\begin{enumerate}
\item \textbf{Free running fails almost everywhere, mainly by overspending.} It
  is outside on every run for 24 of 27 statistics. Customers spend 349 pulls on
  average over the holdout against 110 in the data, choose NotBuy on 9\% of
  occasions against 28\%, and their longest NotBuy run is 3.9 occasions against
  10.9. The worst single statistic is the share of Buy10 on the weapon banner,
  179 times its band. The satiation signature at one copy is 58 times its band,
  far beyond the lottery's own error (8 times, item 5).
\item \textbf{Scheduled sampling halves the damage and keeps the pattern.} At
  $p = 0.35$, 15 of 27 statistics are outside on every run. Spend is still
  double the real level (219 pulls), and the Buy10 weapon share is 72 times its
  band. It does repair the banner switch rate and spend concentration, both
  inside on every run.
\item \textbf{The Markov chain gets the levels right and misses longer memory.}
  Spend, the mean buying run, persistence and escalation are inside on every
  run. What fails is everything that needs more than one step of history: the
  purchase autocorrelation at lags 2--5 (30 to 42 times the band) and the
  longest NotBuy run (16 times). 20 of 27 are outside on every run.
\item \textbf{Independent draws fail everything that depends on order}, as
  designed: the switch rate is 79 times its band and the lag-1 autocorrelation
  101 times. 25 of 27 are outside.
\item \textbf{Arm C passes 26 of 27.} Its one failure is the satiation signature
  at one copy, 7.7 times the band. With real decisions, the simulated lottery
  leaves 68\% of customers without a defined value, against 50\% in the data,
  and those who have one buy again at 0.26 against 0.12. The lottery hands
  featured copies to different customers, or at different moments, from those
  the data record. The Markov chain (7.4) and independent draws (6.9) fail this
  statistic by the same amount, so for them it is the lottery's error, not
  theirs.
\item \textbf{Teacher forcing misses run lengths and switching.} 16 of 27 are
  outside on every run. The worst are the lag-1 autocorrelation (17 times), the
  switch rate (16), the Buy10 figure-B share (9.4), the longest NotBuy run
  (9.2) and the Buy10 regular share (8.6). Spend is 153 pulls against 110. This
  matches the joint classifier, which separates teacher forcing from real data
  mainly on run length and switching.
\end{enumerate}

\section{The 27 tables}

Each table gives, for every arm, the mean $\pm$ s.d.\ of the energy over its
three runs, the mean $z$ against the split-half null, the number of runs inside
the band, the generated statistic's average in its natural units (the real
average is in the caption), and the share of generated histories for which the
statistic is undefined.

\bigskip
% Generated by scripts/r36i_per_statistic.py -- do not edit by hand.
% Source: r36i_functionals.npz; band n_boot 120, seed 3.

\noindent\begin{minipage}{\linewidth}\centering\small
\captionof{table}{\textbf{1. Spend}: total pulls bought over the holdout (Buy1 $=1$, Buy10 $=10$). Real-vs-real band $q_{0.95} = 0.023$; real mean 110.3, undefined for 0\% of customers.}
\begin{tabular}{@{}lrrcrr@{}}
\toprule
 & energy & $z$ & inside & gen.\ mean & undefined\\
\midrule
\multicolumn{6}{@{}l}{\emph{Generators a manager could use}}\\
\quad ProductGPT, free running & $+0.649$ $\pm$ 0.061 & 64.9 & 0/3 & 349.4 & 0\% \\
\quad scheduled sampling, $p = 0.15$ & $+0.298$ $\pm$ 0.106 & 29.9 & 0/3 & 252.9 & 0\% \\
\quad scheduled sampling, $p = 0.35$ & $+0.199$ $\pm$ 0.041 & 19.9 & 0/3 & 218.7 & 0\% \\
\quad first-order Markov chain & $+0.004$ $\pm$ 0.003 & 0.5 & 3/3 & 115.4 & 0\% \\
\quad independent draws & $+0.008$ $\pm$ 0.002 & 0.9 & 3/3 & 116.5 & 0\% \\
\midrule
\multicolumn{6}{@{}l}{\emph{Diagnostic controls: condition on the realised holdout}}\\
\quad arm C: real decisions, lottery re-drawn & $-0.004$ $\pm$ 0.000 & -0.4 & 3/3 & 110.3 & 0\% \\
\quad teacher forcing & $+0.061$ $\pm$ 0.019 & 6.2 & 0/3 & 152.7 & 0\% \\
\bottomrule
\end{tabular}
\end{minipage}\par\bigskip

\noindent\begin{minipage}{\linewidth}\centering\small
\captionof{table}{\textbf{2. Purchase incidence}: share of holdout occasions with a purchase. Real-vs-real band $q_{0.95} = 0.034$; real mean 0.303, undefined for 0\% of customers.}
\begin{tabular}{@{}lrrcrr@{}}
\toprule
 & energy & $z$ & inside & gen.\ mean & undefined\\
\midrule
\multicolumn{6}{@{}l}{\emph{Generators a manager could use}}\\
\quad ProductGPT, free running & $+0.418$ $\pm$ 0.096 & 27.8 & 0/3 & 0.490 & 0\% \\
\quad scheduled sampling, $p = 0.15$ & $+0.203$ $\pm$ 0.091 & 13.5 & 0/3 & 0.431 & 0\% \\
\quad scheduled sampling, $p = 0.35$ & $+0.142$ $\pm$ 0.121 & 9.5 & 1/3 & 0.395 & 0\% \\
\quad first-order Markov chain & $+0.063$ $\pm$ 0.007 & 4.2 & 0/3 & 0.361 & 0\% \\
\quad independent draws & $+0.092$ $\pm$ 0.005 & 6.1 & 0/3 & 0.365 & 0\% \\
\midrule
\multicolumn{6}{@{}l}{\emph{Diagnostic controls: condition on the realised holdout}}\\
\quad arm C: real decisions, lottery re-drawn & $-0.004$ $\pm$ 0.000 & -0.3 & 3/3 & 0.303 & 0\% \\
\quad teacher forcing & $+0.031$ $\pm$ 0.016 & 2.1 & 2/3 & 0.353 & 0\% \\
\bottomrule
\end{tabular}
\end{minipage}\par\bigskip

\noindent\begin{minipage}{\linewidth}\centering\small
\captionof{table}{\textbf{3. Share of Buy1, regular banner}: share of holdout occasions with this decision. Real-vs-real band $q_{0.95} = 0.024$; real mean 0.088, undefined for 0\% of customers.}
\begin{tabular}{@{}lrrcrr@{}}
\toprule
 & energy & $z$ & inside & gen.\ mean & undefined\\
\midrule
\multicolumn{6}{@{}l}{\emph{Generators a manager could use}}\\
\quad ProductGPT, free running & $+0.123$ $\pm$ 0.023 & 7.4 & 0/3 & 0.050 & 0\% \\
\quad scheduled sampling, $p = 0.15$ & $+0.072$ $\pm$ 0.019 & 4.3 & 0/3 & 0.059 & 0\% \\
\quad scheduled sampling, $p = 0.35$ & $+0.085$ $\pm$ 0.021 & 5.1 & 0/3 & 0.057 & 0\% \\
\quad first-order Markov chain & $+0.081$ $\pm$ 0.009 & 4.8 & 0/3 & 0.119 & 0\% \\
\quad independent draws & $+0.112$ $\pm$ 0.009 & 6.7 & 0/3 & 0.119 & 0\% \\
\midrule
\multicolumn{6}{@{}l}{\emph{Diagnostic controls: condition on the realised holdout}}\\
\quad arm C: real decisions, lottery re-drawn & $-0.005$ $\pm$ 0.000 & -0.3 & 3/3 & 0.088 & 0\% \\
\quad teacher forcing & $+0.000$ $\pm$ 0.004 & -0.0 & 3/3 & 0.087 & 0\% \\
\bottomrule
\end{tabular}
\end{minipage}\par\bigskip

\noindent\begin{minipage}{\linewidth}\centering\small
\captionof{table}{\textbf{4. Share of Buy10, regular banner}: share of holdout occasions with this decision. Real-vs-real band $q_{0.95} = 0.015$; real mean 0.003, undefined for 0\% of customers.}
\begin{tabular}{@{}lrrcrr@{}}
\toprule
 & energy & $z$ & inside & gen.\ mean & undefined\\
\midrule
\multicolumn{6}{@{}l}{\emph{Generators a manager could use}}\\
\quad ProductGPT, free running & $+0.290$ $\pm$ 0.035 & 39.8 & 0/3 & 0.017 & 0\% \\
\quad scheduled sampling, $p = 0.15$ & $+0.293$ $\pm$ 0.068 & 40.3 & 0/3 & 0.013 & 0\% \\
\quad scheduled sampling, $p = 0.35$ & $+0.113$ $\pm$ 0.026 & 15.5 & 0/3 & 0.009 & 0\% \\
\quad first-order Markov chain & $+0.139$ $\pm$ 0.020 & 19.2 & 0/3 & 0.006 & 0\% \\
\quad independent draws & $+0.179$ $\pm$ 0.009 & 24.6 & 0/3 & 0.006 & 0\% \\
\midrule
\multicolumn{6}{@{}l}{\emph{Diagnostic controls: condition on the realised holdout}}\\
\quad arm C: real decisions, lottery re-drawn & $-0.002$ $\pm$ 0.000 & -0.2 & 3/3 & 0.003 & 0\% \\
\quad teacher forcing & $+0.125$ $\pm$ 0.041 & 17.2 & 0/3 & 0.009 & 0\% \\
\bottomrule
\end{tabular}
\end{minipage}\par\bigskip

\noindent\begin{minipage}{\linewidth}\centering\small
\captionof{table}{\textbf{5. Share of Buy1, figure A}: share of holdout occasions with this decision. Real-vs-real band $q_{0.95} = 0.032$; real mean 0.104, undefined for 0\% of customers.}
\begin{tabular}{@{}lrrcrr@{}}
\toprule
 & energy & $z$ & inside & gen.\ mean & undefined\\
\midrule
\multicolumn{6}{@{}l}{\emph{Generators a manager could use}}\\
\quad ProductGPT, free running & $+0.112$ $\pm$ 0.055 & 6.9 & 0/3 & 0.047 & 0\% \\
\quad scheduled sampling, $p = 0.15$ & $+0.105$ $\pm$ 0.025 & 6.5 & 0/3 & 0.046 & 0\% \\
\quad scheduled sampling, $p = 0.35$ & $+0.075$ $\pm$ 0.042 & 4.6 & 1/3 & 0.055 & 0\% \\
\quad first-order Markov chain & $+0.067$ $\pm$ 0.002 & 4.1 & 0/3 & 0.107 & 0\% \\
\quad independent draws & $+0.128$ $\pm$ 0.003 & 7.9 & 0/3 & 0.109 & 0\% \\
\midrule
\multicolumn{6}{@{}l}{\emph{Diagnostic controls: condition on the realised holdout}}\\
\quad arm C: real decisions, lottery re-drawn & $-0.002$ $\pm$ 0.000 & -0.2 & 3/3 & 0.104 & 0\% \\
\quad teacher forcing & $-0.000$ $\pm$ 0.000 & -0.1 & 3/3 & 0.101 & 0\% \\
\bottomrule
\end{tabular}
\end{minipage}\par\bigskip

\noindent\begin{minipage}{\linewidth}\centering\small
\captionof{table}{\textbf{6. Share of Buy10, figure A}: share of holdout occasions with this decision. Real-vs-real band $q_{0.95} = 0.017$; real mean 0.045, undefined for 0\% of customers.}
\begin{tabular}{@{}lrrcrr@{}}
\toprule
 & energy & $z$ & inside & gen.\ mean & undefined\\
\midrule
\multicolumn{6}{@{}l}{\emph{Generators a manager could use}}\\
\quad ProductGPT, free running & $+0.022$ $\pm$ 0.007 & 2.3 & 1/3 & 0.059 & 0\% \\
\quad scheduled sampling, $p = 0.15$ & $+0.005$ $\pm$ 0.004 & 0.6 & 3/3 & 0.045 & 0\% \\
\quad scheduled sampling, $p = 0.35$ & $+0.010$ $\pm$ 0.013 & 1.1 & 2/3 & 0.042 & 0\% \\
\quad first-order Markov chain & $+0.018$ $\pm$ 0.005 & 1.8 & 1/3 & 0.031 & 0\% \\
\quad independent draws & $+0.052$ $\pm$ 0.003 & 5.1 & 0/3 & 0.032 & 0\% \\
\midrule
\multicolumn{6}{@{}l}{\emph{Diagnostic controls: condition on the realised holdout}}\\
\quad arm C: real decisions, lottery re-drawn & $-0.004$ $\pm$ 0.000 & -0.3 & 3/3 & 0.045 & 0\% \\
\quad teacher forcing & $+0.016$ $\pm$ 0.005 & 1.6 & 2/3 & 0.056 & 0\% \\
\bottomrule
\end{tabular}
\end{minipage}\par\bigskip

\noindent\begin{minipage}{\linewidth}\centering\small
\captionof{table}{\textbf{7. Share of Buy1, figure B}: share of holdout occasions with this decision. Real-vs-real band $q_{0.95} = 0.017$; real mean 0.030, undefined for 0\% of customers.}
\begin{tabular}{@{}lrrcrr@{}}
\toprule
 & energy & $z$ & inside & gen.\ mean & undefined\\
\midrule
\multicolumn{6}{@{}l}{\emph{Generators a manager could use}}\\
\quad ProductGPT, free running & $+0.078$ $\pm$ 0.090 & 9.4 & 0/3 & 0.063 & 0\% \\
\quad scheduled sampling, $p = 0.15$ & $+0.089$ $\pm$ 0.035 & 10.7 & 0/3 & 0.068 & 0\% \\
\quad scheduled sampling, $p = 0.35$ & $+0.130$ $\pm$ 0.199 & 15.7 & 1/3 & 0.057 & 0\% \\
\quad first-order Markov chain & $+0.203$ $\pm$ 0.013 & 24.3 & 0/3 & 0.056 & 0\% \\
\quad independent draws & $+0.332$ $\pm$ 0.004 & 39.8 & 0/3 & 0.057 & 0\% \\
\midrule
\multicolumn{6}{@{}l}{\emph{Diagnostic controls: condition on the realised holdout}}\\
\quad arm C: real decisions, lottery re-drawn & $-0.003$ $\pm$ 0.000 & -0.3 & 3/3 & 0.030 & 0\% \\
\quad teacher forcing & $+0.007$ $\pm$ 0.003 & 0.9 & 3/3 & 0.038 & 0\% \\
\bottomrule
\end{tabular}
\end{minipage}\par\bigskip

\noindent\begin{minipage}{\linewidth}\centering\small
\captionof{table}{\textbf{8. Share of Buy10, figure B}: share of holdout occasions with this decision. Real-vs-real band $q_{0.95} = 0.018$; real mean 0.008, undefined for 0\% of customers.}
\begin{tabular}{@{}lrrcrr@{}}
\toprule
 & energy & $z$ & inside & gen.\ mean & undefined\\
\midrule
\multicolumn{6}{@{}l}{\emph{Generators a manager could use}}\\
\quad ProductGPT, free running & $+0.931$ $\pm$ 0.533 & 109.5 & 0/3 & 0.068 & 0\% \\
\quad scheduled sampling, $p = 0.15$ & $+0.632$ $\pm$ 0.222 & 74.3 & 0/3 & 0.054 & 0\% \\
\quad scheduled sampling, $p = 0.35$ & $+0.323$ $\pm$ 0.403 & 38.0 & 0/3 & 0.034 & 0\% \\
\quad first-order Markov chain & $+0.069$ $\pm$ 0.003 & 8.1 & 0/3 & 0.014 & 0\% \\
\quad independent draws & $+0.162$ $\pm$ 0.011 & 19.0 & 0/3 & 0.014 & 0\% \\
\midrule
\multicolumn{6}{@{}l}{\emph{Diagnostic controls: condition on the realised holdout}}\\
\quad arm C: real decisions, lottery re-drawn & $-0.003$ $\pm$ 0.000 & -0.4 & 3/3 & 0.008 & 0\% \\
\quad teacher forcing & $+0.166$ $\pm$ 0.070 & 19.4 & 0/3 & 0.021 & 0\% \\
\bottomrule
\end{tabular}
\end{minipage}\par\bigskip

\noindent\begin{minipage}{\linewidth}\centering\small
\captionof{table}{\textbf{9. Share of Buy1, weapon banner}: share of holdout occasions with this decision. Real-vs-real band $q_{0.95} = 0.016$; real mean 0.019, undefined for 0\% of customers.}
\begin{tabular}{@{}lrrcrr@{}}
\toprule
 & energy & $z$ & inside & gen.\ mean & undefined\\
\midrule
\multicolumn{6}{@{}l}{\emph{Generators a manager could use}}\\
\quad ProductGPT, free running & $+0.348$ $\pm$ 0.164 & 44.1 & 0/3 & 0.087 & 0\% \\
\quad scheduled sampling, $p = 0.15$ & $+0.293$ $\pm$ 0.065 & 37.2 & 0/3 & 0.082 & 0\% \\
\quad scheduled sampling, $p = 0.35$ & $+0.283$ $\pm$ 0.123 & 35.9 & 0/3 & 0.080 & 0\% \\
\quad first-order Markov chain & $+0.042$ $\pm$ 0.001 & 5.2 & 0/3 & 0.020 & 0\% \\
\quad independent draws & $+0.117$ $\pm$ 0.003 & 14.8 & 0/3 & 0.020 & 0\% \\
\midrule
\multicolumn{6}{@{}l}{\emph{Diagnostic controls: condition on the realised holdout}}\\
\quad arm C: real decisions, lottery re-drawn & $-0.002$ $\pm$ 0.000 & -0.3 & 3/3 & 0.019 & 0\% \\
\quad teacher forcing & $+0.016$ $\pm$ 0.002 & 2.0 & 2/3 & 0.028 & 0\% \\
\bottomrule
\end{tabular}
\end{minipage}\par\bigskip

\noindent\begin{minipage}{\linewidth}\centering\small
\captionof{table}{\textbf{10. Share of Buy10, weapon banner}: share of holdout occasions with this decision. Real-vs-real band $q_{0.95} = 0.012$; real mean 0.006, undefined for 0\% of customers.}
\begin{tabular}{@{}lrrcrr@{}}
\toprule
 & energy & $z$ & inside & gen.\ mean & undefined\\
\midrule
\multicolumn{6}{@{}l}{\emph{Generators a manager could use}}\\
\quad ProductGPT, free running & $+2.190$ $\pm$ 0.228 & 361.0 & 0/3 & 0.099 & 0\% \\
\quad scheduled sampling, $p = 0.15$ & $+0.943$ $\pm$ 0.321 & 155.5 & 0/3 & 0.063 & 0\% \\
\quad scheduled sampling, $p = 0.35$ & $+0.878$ $\pm$ 0.197 & 144.8 & 0/3 & 0.061 & 0\% \\
\quad first-order Markov chain & $+0.013$ $\pm$ 0.006 & 2.1 & 1/3 & 0.008 & 0\% \\
\quad independent draws & $+0.087$ $\pm$ 0.008 & 14.3 & 0/3 & 0.008 & 0\% \\
\midrule
\multicolumn{6}{@{}l}{\emph{Diagnostic controls: condition on the realised holdout}}\\
\quad arm C: real decisions, lottery re-drawn & $-0.002$ $\pm$ 0.000 & -0.2 & 3/3 & 0.006 & 0\% \\
\quad teacher forcing & $+0.056$ $\pm$ 0.007 & 9.3 & 0/3 & 0.013 & 0\% \\
\bottomrule
\end{tabular}
\end{minipage}\par\bigskip

\noindent\begin{minipage}{\linewidth}\centering\small
\captionof{table}{\textbf{11. Share of NotBuy}: share of holdout occasions with this decision. Real-vs-real band $q_{0.95} = 0.026$; real mean 0.280, undefined for 0\% of customers.}
\begin{tabular}{@{}lrrcrr@{}}
\toprule
 & energy & $z$ & inside & gen.\ mean & undefined\\
\midrule
\multicolumn{6}{@{}l}{\emph{Generators a manager could use}}\\
\quad ProductGPT, free running & $+0.425$ $\pm$ 0.100 & 36.4 & 0/3 & 0.093 & 0\% \\
\quad scheduled sampling, $p = 0.15$ & $+0.204$ $\pm$ 0.095 & 17.5 & 0/3 & 0.153 & 0\% \\
\quad scheduled sampling, $p = 0.35$ & $+0.142$ $\pm$ 0.122 & 12.2 & 1/3 & 0.188 & 0\% \\
\quad first-order Markov chain & $+0.061$ $\pm$ 0.007 & 5.3 & 0/3 & 0.222 & 0\% \\
\quad independent draws & $+0.091$ $\pm$ 0.005 & 7.8 & 0/3 & 0.218 & 0\% \\
\midrule
\multicolumn{6}{@{}l}{\emph{Diagnostic controls: condition on the realised holdout}}\\
\quad arm C: real decisions, lottery re-drawn & $-0.004$ $\pm$ 0.000 & -0.2 & 3/3 & 0.280 & 0\% \\
\quad teacher forcing & $+0.028$ $\pm$ 0.016 & 2.5 & 1/3 & 0.230 & 0\% \\
\bottomrule
\end{tabular}
\end{minipage}\par\bigskip

\noindent\begin{minipage}{\linewidth}\centering\small
\captionof{table}{\textbf{12. Longest NotBuy run}: longest run of consecutive NotBuy occasions. Real-vs-real band $q_{0.95} = 0.025$; real mean 10.9, undefined for 0\% of customers.}
\begin{tabular}{@{}lrrcrr@{}}
\toprule
 & energy & $z$ & inside & gen.\ mean & undefined\\
\midrule
\multicolumn{6}{@{}l}{\emph{Generators a manager could use}}\\
\quad ProductGPT, free running & $+0.517$ $\pm$ 0.132 & 39.8 & 0/3 & 3.918 & 0\% \\
\quad scheduled sampling, $p = 0.15$ & $+0.295$ $\pm$ 0.096 & 22.8 & 0/3 & 7.161 & 0\% \\
\quad scheduled sampling, $p = 0.35$ & $+0.241$ $\pm$ 0.088 & 18.6 & 0/3 & 9.147 & 0\% \\
\quad first-order Markov chain & $+0.398$ $\pm$ 0.009 & 30.7 & 0/3 & 4.775 & 0\% \\
\quad independent draws & $+0.695$ $\pm$ 0.005 & 53.5 & 0/3 & 2.575 & 0\% \\
\midrule
\multicolumn{6}{@{}l}{\emph{Diagnostic controls: condition on the realised holdout}}\\
\quad arm C: real decisions, lottery re-drawn & $-0.003$ $\pm$ 0.000 & -0.1 & 3/3 & 10.9 & 0\% \\
\quad teacher forcing & $+0.232$ $\pm$ 0.062 & 17.9 & 0/3 & 6.137 & 0\% \\
\bottomrule
\end{tabular}
\end{minipage}\par\bigskip

\noindent\begin{minipage}{\linewidth}\centering\small
\captionof{table}{\textbf{13. Mean buying run}: mean length of runs of consecutive purchases. Real-vs-real band $q_{0.95} = 0.023$; real mean 3.431, undefined for 0\% of customers.}
\begin{tabular}{@{}lrrcrr@{}}
\toprule
 & energy & $z$ & inside & gen.\ mean & undefined\\
\midrule
\multicolumn{6}{@{}l}{\emph{Generators a manager could use}}\\
\quad ProductGPT, free running & $+0.671$ $\pm$ 0.129 & 60.3 & 0/3 & 11.9 & 0\% \\
\quad scheduled sampling, $p = 0.15$ & $+0.163$ $\pm$ 0.098 & 14.7 & 0/3 & 6.463 & 0\% \\
\quad scheduled sampling, $p = 0.35$ & $+0.145$ $\pm$ 0.058 & 13.1 & 0/3 & 6.299 & 0\% \\
\quad first-order Markov chain & $+0.020$ $\pm$ 0.002 & 1.8 & 3/3 & 2.765 & 0\% \\
\quad independent draws & $+0.233$ $\pm$ 0.002 & 21.0 & 0/3 & 1.559 & 0\% \\
\midrule
\multicolumn{6}{@{}l}{\emph{Diagnostic controls: condition on the realised holdout}}\\
\quad arm C: real decisions, lottery re-drawn & $-0.003$ $\pm$ 0.000 & -0.3 & 3/3 & 3.431 & 0\% \\
\quad teacher forcing & $+0.018$ $\pm$ 0.005 & 1.6 & 2/3 & 2.713 & 0\% \\
\bottomrule
\end{tabular}
\end{minipage}\par\bigskip

\noindent\begin{minipage}{\linewidth}\centering\small
\captionof{table}{\textbf{14. Longest buying run}: longest run of consecutive purchases. Real-vs-real band $q_{0.95} = 0.028$; real mean 11.9, undefined for 0\% of customers.}
\begin{tabular}{@{}lrrcrr@{}}
\toprule
 & energy & $z$ & inside & gen.\ mean & undefined\\
\midrule
\multicolumn{6}{@{}l}{\emph{Generators a manager could use}}\\
\quad ProductGPT, free running & $+0.419$ $\pm$ 0.079 & 31.7 & 0/3 & 28.8 & 0\% \\
\quad scheduled sampling, $p = 0.15$ & $+0.129$ $\pm$ 0.069 & 9.7 & 0/3 & 19.8 & 0\% \\
\quad scheduled sampling, $p = 0.35$ & $+0.117$ $\pm$ 0.040 & 8.8 & 0/3 & 19.1 & 0\% \\
\quad first-order Markov chain & $+0.031$ $\pm$ 0.001 & 2.3 & 0/3 & 8.540 & 0\% \\
\quad independent draws & $+0.188$ $\pm$ 0.003 & 14.2 & 0/3 & 5.106 & 0\% \\
\midrule
\multicolumn{6}{@{}l}{\emph{Diagnostic controls: condition on the realised holdout}}\\
\quad arm C: real decisions, lottery re-drawn & $-0.004$ $\pm$ 0.000 & -0.3 & 3/3 & 11.9 & 0\% \\
\quad teacher forcing & $-0.002$ $\pm$ 0.001 & -0.2 & 3/3 & 11.6 & 0\% \\
\bottomrule
\end{tabular}
\end{minipage}\par\bigskip

\noindent\begin{minipage}{\linewidth}\centering\small
\captionof{table}{\textbf{15. Banner switch rate}: share of consecutive purchases that change banner. Real-vs-real band $q_{0.95} = 0.025$; real mean 0.126, undefined for 0\% of customers.}
\begin{tabular}{@{}lrrcrr@{}}
\toprule
 & energy & $z$ & inside & gen.\ mean & undefined\\
\midrule
\multicolumn{6}{@{}l}{\emph{Generators a manager could use}}\\
\quad ProductGPT, free running & $+0.030$ $\pm$ 0.010 & 2.5 & 1/3 & 0.102 & 0\% \\
\quad scheduled sampling, $p = 0.15$ & $-0.001$ $\pm$ 0.001 & -0.0 & 3/3 & 0.127 & 0\% \\
\quad scheduled sampling, $p = 0.35$ & $+0.009$ $\pm$ 0.013 & 0.8 & 3/3 & 0.125 & 0\% \\
\quad first-order Markov chain & $+0.178$ $\pm$ 0.013 & 14.6 & 0/3 & 0.182 & 0\% \\
\quad independent draws & $+2.005$ $\pm$ 0.012 & 164.4 & 0/3 & 0.403 & 0\% \\
\midrule
\multicolumn{6}{@{}l}{\emph{Diagnostic controls: condition on the realised holdout}}\\
\quad arm C: real decisions, lottery re-drawn & $-0.004$ $\pm$ 0.000 & -0.3 & 3/3 & 0.126 & 0\% \\
\quad teacher forcing & $+0.413$ $\pm$ 0.051 & 34.0 & 0/3 & 0.233 & 0\% \\
\bottomrule
\end{tabular}
\end{minipage}\par\bigskip

\noindent\begin{minipage}{\linewidth}\centering\small
\captionof{table}{\textbf{16. Spend concentration}: Herfindahl index of spend across the four banners. Real-vs-real band $q_{0.95} = 0.024$; real mean 0.332, undefined for 0\% of customers.}
\begin{tabular}{@{}lrrcrr@{}}
\toprule
 & energy & $z$ & inside & gen.\ mean & undefined\\
\midrule
\multicolumn{6}{@{}l}{\emph{Generators a manager could use}}\\
\quad ProductGPT, free running & $+0.025$ $\pm$ 0.004 & 1.7 & 2/3 & 0.284 & 0\% \\
\quad scheduled sampling, $p = 0.15$ & $+0.027$ $\pm$ 0.002 & 1.8 & 0/3 & 0.283 & 0\% \\
\quad scheduled sampling, $p = 0.35$ & $+0.012$ $\pm$ 0.009 & 0.9 & 3/3 & 0.301 & 0\% \\
\quad first-order Markov chain & $+0.108$ $\pm$ 0.006 & 7.3 & 0/3 & 0.234 & 0\% \\
\quad independent draws & $+0.188$ $\pm$ 0.011 & 12.7 & 0/3 & 0.207 & 0\% \\
\midrule
\multicolumn{6}{@{}l}{\emph{Diagnostic controls: condition on the realised holdout}}\\
\quad arm C: real decisions, lottery re-drawn & $-0.003$ $\pm$ 0.000 & -0.1 & 3/3 & 0.332 & 0\% \\
\quad teacher forcing & $+0.047$ $\pm$ 0.010 & 3.2 & 0/3 & 0.265 & 0\% \\
\bottomrule
\end{tabular}
\end{minipage}\par\bigskip

\noindent\begin{minipage}{\linewidth}\centering\small
\captionof{table}{\textbf{17. Autocorrelation, lag 1}: autocorrelation of the purchase indicator at lag 1. Real-vs-real band $q_{0.95} = 0.013$; real mean 0.546, undefined for 42\% of customers.}
\begin{tabular}{@{}lrrcrr@{}}
\toprule
 & energy & $z$ & inside & gen.\ mean & undefined\\
\midrule
\multicolumn{6}{@{}l}{\emph{Generators a manager could use}}\\
\quad ProductGPT, free running & $+0.302$ $\pm$ 0.140 & 47.6 & 0/3 & 0.267 & 45\% \\
\quad scheduled sampling, $p = 0.15$ & $+0.160$ $\pm$ 0.074 & 25.2 & 0/3 & 0.348 & 43\% \\
\quad scheduled sampling, $p = 0.35$ & $+0.110$ $\pm$ 0.095 & 17.3 & 1/3 & 0.423 & 43\% \\
\quad first-order Markov chain & $+0.152$ $\pm$ 0.001 & 23.9 & 0/3 & 0.419 & 42\% \\
\quad independent draws & $+1.316$ $\pm$ 0.020 & 206.7 & 0/3 & -0.012 & 42\% \\
\midrule
\multicolumn{6}{@{}l}{\emph{Diagnostic controls: condition on the realised holdout}}\\
\quad arm C: real decisions, lottery re-drawn & $-0.002$ $\pm$ 0.000 & -0.2 & 3/3 & 0.546 & 42\% \\
\quad teacher forcing & $+0.221$ $\pm$ 0.040 & 34.8 & 0/3 & 0.346 & 42\% \\
\bottomrule
\end{tabular}
\end{minipage}\par\bigskip

\noindent\begin{minipage}{\linewidth}\centering\small
\captionof{table}{\textbf{18. Autocorrelation, lag 2}: autocorrelation of the purchase indicator at lag 2. Real-vs-real band $q_{0.95} = 0.015$; real mean 0.459, undefined for 42\% of customers.}
\begin{tabular}{@{}lrrcrr@{}}
\toprule
 & energy & $z$ & inside & gen.\ mean & undefined\\
\midrule
\multicolumn{6}{@{}l}{\emph{Generators a manager could use}}\\
\quad ProductGPT, free running & $+0.283$ $\pm$ 0.123 & 40.0 & 0/3 & 0.214 & 46\% \\
\quad scheduled sampling, $p = 0.15$ & $+0.167$ $\pm$ 0.080 & 23.6 & 0/3 & 0.278 & 43\% \\
\quad scheduled sampling, $p = 0.35$ & $+0.112$ $\pm$ 0.092 & 15.8 & 1/3 & 0.338 & 43\% \\
\quad first-order Markov chain & $+0.469$ $\pm$ 0.002 & 66.0 & 0/3 & 0.177 & 42\% \\
\quad independent draws & $+1.137$ $\pm$ 0.010 & 160.2 & 0/3 & -0.014 & 42\% \\
\midrule
\multicolumn{6}{@{}l}{\emph{Diagnostic controls: condition on the realised holdout}}\\
\quad arm C: real decisions, lottery re-drawn & $-0.002$ $\pm$ 0.000 & -0.2 & 3/3 & 0.459 & 42\% \\
\quad teacher forcing & $+0.129$ $\pm$ 0.034 & 18.2 & 0/3 & 0.310 & 42\% \\
\bottomrule
\end{tabular}
\end{minipage}\par\bigskip

\noindent\begin{minipage}{\linewidth}\centering\small
\captionof{table}{\textbf{19. Autocorrelation, lag 3}: autocorrelation of the purchase indicator at lag 3. Real-vs-real band $q_{0.95} = 0.013$; real mean 0.386, undefined for 42\% of customers.}
\begin{tabular}{@{}lrrcrr@{}}
\toprule
 & energy & $z$ & inside & gen.\ mean & undefined\\
\midrule
\multicolumn{6}{@{}l}{\emph{Generators a manager could use}}\\
\quad ProductGPT, free running & $+0.225$ $\pm$ 0.093 & 32.4 & 0/3 & 0.172 & 46\% \\
\quad scheduled sampling, $p = 0.15$ & $+0.132$ $\pm$ 0.065 & 19.0 & 0/3 & 0.228 & 43\% \\
\quad scheduled sampling, $p = 0.35$ & $+0.091$ $\pm$ 0.073 & 13.1 & 1/3 & 0.280 & 43\% \\
\quad first-order Markov chain & $+0.556$ $\pm$ 0.005 & 80.0 & 0/3 & 0.067 & 42\% \\
\quad independent draws & $+0.838$ $\pm$ 0.009 & 120.7 & 0/3 & -0.011 & 42\% \\
\midrule
\multicolumn{6}{@{}l}{\emph{Diagnostic controls: condition on the realised holdout}}\\
\quad arm C: real decisions, lottery re-drawn & $-0.003$ $\pm$ 0.000 & -0.4 & 3/3 & 0.386 & 42\% \\
\quad teacher forcing & $+0.079$ $\pm$ 0.030 & 11.3 & 0/3 & 0.269 & 42\% \\
\bottomrule
\end{tabular}
\end{minipage}\par\bigskip

\noindent\begin{minipage}{\linewidth}\centering\small
\captionof{table}{\textbf{20. Autocorrelation, lag 4}: autocorrelation of the purchase indicator at lag 4. Real-vs-real band $q_{0.95} = 0.014$; real mean 0.338, undefined for 42\% of customers.}
\begin{tabular}{@{}lrrcrr@{}}
\toprule
 & energy & $z$ & inside & gen.\ mean & undefined\\
\midrule
\multicolumn{6}{@{}l}{\emph{Generators a manager could use}}\\
\quad ProductGPT, free running & $+0.207$ $\pm$ 0.085 & 30.7 & 0/3 & 0.141 & 46\% \\
\quad scheduled sampling, $p = 0.15$ & $+0.122$ $\pm$ 0.058 & 18.2 & 0/3 & 0.194 & 43\% \\
\quad scheduled sampling, $p = 0.35$ & $+0.085$ $\pm$ 0.063 & 12.6 & 0/3 & 0.240 & 44\% \\
\quad first-order Markov chain & $+0.564$ $\pm$ 0.007 & 83.7 & 0/3 & 0.018 & 42\% \\
\quad independent draws & $+0.686$ $\pm$ 0.022 & 101.7 & 0/3 & -0.009 & 42\% \\
\midrule
\multicolumn{6}{@{}l}{\emph{Diagnostic controls: condition on the realised holdout}}\\
\quad arm C: real decisions, lottery re-drawn & $-0.003$ $\pm$ 0.000 & -0.4 & 3/3 & 0.338 & 42\% \\
\quad teacher forcing & $+0.067$ $\pm$ 0.024 & 10.0 & 0/3 & 0.232 & 42\% \\
\bottomrule
\end{tabular}
\end{minipage}\par\bigskip

\noindent\begin{minipage}{\linewidth}\centering\small
\captionof{table}{\textbf{21. Autocorrelation, lag 5}: autocorrelation of the purchase indicator at lag 5. Real-vs-real band $q_{0.95} = 0.017$; real mean 0.285, undefined for 42\% of customers.}
\begin{tabular}{@{}lrrcrr@{}}
\toprule
 & energy & $z$ & inside & gen.\ mean & undefined\\
\midrule
\multicolumn{6}{@{}l}{\emph{Generators a manager could use}}\\
\quad ProductGPT, free running & $+0.165$ $\pm$ 0.074 & 22.0 & 0/3 & 0.120 & 46\% \\
\quad scheduled sampling, $p = 0.15$ & $+0.098$ $\pm$ 0.043 & 13.0 & 0/3 & 0.166 & 43\% \\
\quad scheduled sampling, $p = 0.35$ & $+0.070$ $\pm$ 0.044 & 9.3 & 0/3 & 0.209 & 44\% \\
\quad first-order Markov chain & $+0.491$ $\pm$ 0.012 & 65.5 & 0/3 & -0.005 & 42\% \\
\quad independent draws & $+0.537$ $\pm$ 0.015 & 71.7 & 0/3 & -0.010 & 42\% \\
\midrule
\multicolumn{6}{@{}l}{\emph{Diagnostic controls: condition on the realised holdout}}\\
\quad arm C: real decisions, lottery re-drawn & $-0.002$ $\pm$ 0.000 & -0.3 & 3/3 & 0.285 & 42\% \\
\quad teacher forcing & $+0.037$ $\pm$ 0.015 & 4.9 & 0/3 & 0.207 & 42\% \\
\bottomrule
\end{tabular}
\end{minipage}\par\bigskip

\noindent\begin{minipage}{\linewidth}\centering\small
\captionof{table}{\textbf{22. Escalation}: $P(\text{Buy10 at } t \mid \text{Buy10 at } t-1)$. Real-vs-real band $q_{0.95} = 0.017$; real mean 0.535, undefined for 46\% of customers.}
\begin{tabular}{@{}lrrcrr@{}}
\toprule
 & energy & $z$ & inside & gen.\ mean & undefined\\
\midrule
\multicolumn{6}{@{}l}{\emph{Generators a manager could use}}\\
\quad ProductGPT, free running & $+0.127$ $\pm$ 0.002 & 15.2 & 0/3 & 0.716 & 43\% \\
\quad scheduled sampling, $p = 0.15$ & $+0.024$ $\pm$ 0.017 & 2.8 & 2/3 & 0.610 & 43\% \\
\quad scheduled sampling, $p = 0.35$ & $+0.017$ $\pm$ 0.021 & 2.0 & 2/3 & 0.591 & 45\% \\
\quad first-order Markov chain & $+0.006$ $\pm$ 0.001 & 0.6 & 3/3 & 0.553 & 44\% \\
\quad independent draws & $+0.735$ $\pm$ 0.008 & 88.0 & 0/3 & 0.093 & 42\% \\
\midrule
\multicolumn{6}{@{}l}{\emph{Diagnostic controls: condition on the realised holdout}}\\
\quad arm C: real decisions, lottery re-drawn & $-0.003$ $\pm$ 0.000 & -0.5 & 3/3 & 0.535 & 46\% \\
\quad teacher forcing & $+0.100$ $\pm$ 0.010 & 11.9 & 0/3 & 0.379 & 43\% \\
\bottomrule
\end{tabular}
\end{minipage}\par\bigskip

\noindent\begin{minipage}{\linewidth}\centering\small
\captionof{table}{\textbf{23. Persistence}: $P(\text{purchase at } t \mid \text{purchase at } t-1)$. Real-vs-real band $q_{0.95} = 0.014$; real mean 0.787, undefined for 42\% of customers.}
\begin{tabular}{@{}lrrcrr@{}}
\toprule
 & energy & $z$ & inside & gen.\ mean & undefined\\
\midrule
\multicolumn{6}{@{}l}{\emph{Generators a manager could use}}\\
\quad ProductGPT, free running & $+0.315$ $\pm$ 0.069 & 49.0 & 0/3 & 0.901 & 42\% \\
\quad scheduled sampling, $p = 0.15$ & $+0.135$ $\pm$ 0.075 & 20.9 & 0/3 & 0.857 & 42\% \\
\quad scheduled sampling, $p = 0.35$ & $+0.122$ $\pm$ 0.057 & 19.0 & 0/3 & 0.851 & 42\% \\
\quad first-order Markov chain & $+0.011$ $\pm$ 0.001 & 1.7 & 3/3 & 0.781 & 42\% \\
\quad independent draws & $+0.643$ $\pm$ 0.014 & 100.0 & 0/3 & 0.620 & 42\% \\
\midrule
\multicolumn{6}{@{}l}{\emph{Diagnostic controls: condition on the realised holdout}}\\
\quad arm C: real decisions, lottery re-drawn & $-0.002$ $\pm$ 0.000 & -0.4 & 3/3 & 0.787 & 42\% \\
\quad teacher forcing & $+0.038$ $\pm$ 0.011 & 5.9 & 0/3 & 0.742 & 42\% \\
\bottomrule
\end{tabular}
\end{minipage}\par\bigskip

\noindent\begin{minipage}{\linewidth}\centering\small
\captionof{table}{\textbf{24. Satiation signature, $k = 0$}: $P(\text{buy on a banner} \mid \text{holding 0 copies of its featured product})$, pooled over figure A, figure B and weapon. Real-vs-real band $q_{0.95} = 0.015$; real mean 0.135, undefined for 42\% of customers.}
\begin{tabular}{@{}lrrcrr@{}}
\toprule
 & energy & $z$ & inside & gen.\ mean & undefined\\
\midrule
\multicolumn{6}{@{}l}{\emph{Generators a manager could use}}\\
\quad ProductGPT, free running & $+0.283$ $\pm$ 0.034 & 32.2 & 0/3 & 0.217 & 42\% \\
\quad scheduled sampling, $p = 0.15$ & $+0.139$ $\pm$ 0.076 & 15.8 & 0/3 & 0.189 & 42\% \\
\quad scheduled sampling, $p = 0.35$ & $+0.121$ $\pm$ 0.117 & 13.8 & 1/3 & 0.176 & 42\% \\
\quad first-order Markov chain & $+0.042$ $\pm$ 0.001 & 4.8 & 0/3 & 0.131 & 42\% \\
\quad independent draws & $+0.082$ $\pm$ 0.001 & 9.4 & 0/3 & 0.136 & 42\% \\
\midrule
\multicolumn{6}{@{}l}{\emph{Diagnostic controls: condition on the realised holdout}}\\
\quad arm C: real decisions, lottery re-drawn & $+0.010$ $\pm$ 0.002 & 1.2 & 3/3 & 0.118 & 42\% \\
\quad teacher forcing & $+0.029$ $\pm$ 0.005 & 3.4 & 0/3 & 0.161 & 42\% \\
\bottomrule
\end{tabular}
\end{minipage}\par\bigskip

\noindent\begin{minipage}{\linewidth}\centering\small
\captionof{table}{\textbf{25. Satiation signature, $k = 1$}: $P(\text{buy on a banner} \mid \text{holding 1 copies of its featured product})$, pooled over figure A, figure B and weapon. Real-vs-real band $q_{0.95} = 0.008$; real mean 0.124, undefined for 50\% of customers.}
\begin{tabular}{@{}lrrcrr@{}}
\toprule
 & energy & $z$ & inside & gen.\ mean & undefined\\
\midrule
\multicolumn{6}{@{}l}{\emph{Generators a manager could use}}\\
\quad ProductGPT, free running & $+0.439$ $\pm$ 0.070 & 96.9 & 0/3 & 0.522 & 55\% \\
\quad scheduled sampling, $p = 0.15$ & $+0.280$ $\pm$ 0.039 & 61.9 & 0/3 & 0.470 & 59\% \\
\quad scheduled sampling, $p = 0.35$ & $+0.257$ $\pm$ 0.037 & 56.9 & 0/3 & 0.467 & 60\% \\
\quad first-order Markov chain & $+0.057$ $\pm$ 0.002 & 12.7 & 0/3 & 0.239 & 68\% \\
\quad independent draws & $+0.053$ $\pm$ 0.003 & 11.8 & 0/3 & 0.174 & 68\% \\
\midrule
\multicolumn{6}{@{}l}{\emph{Diagnostic controls: condition on the realised holdout}}\\
\quad arm C: real decisions, lottery re-drawn & $+0.058$ $\pm$ 0.002 & 13.0 & 0/3 & 0.264 & 68\% \\
\quad teacher forcing & $+0.014$ $\pm$ 0.005 & 3.3 & 0/3 & 0.164 & 50\% \\
\bottomrule
\end{tabular}
\end{minipage}\par\bigskip

\noindent\begin{minipage}{\linewidth}\centering\small
\captionof{table}{\textbf{26. Satiation signature, $k = 2$}: $P(\text{buy on a banner} \mid \text{holding 2 copies of its featured product})$, pooled over figure A, figure B and weapon. Real-vs-real band $q_{0.95} = 0.003$; real mean 0.213, undefined for 89\% of customers.}
\begin{tabular}{@{}lrrcrr@{}}
\toprule
 & energy & $z$ & inside & gen.\ mean & undefined\\
\midrule
\multicolumn{6}{@{}l}{\emph{Generators a manager could use}}\\
\quad ProductGPT, free running & $+0.069$ $\pm$ 0.010 & 40.8 & 0/3 & 0.673 & 79\% \\
\quad scheduled sampling, $p = 0.15$ & $+0.023$ $\pm$ 0.007 & 13.6 & 0/3 & 0.609 & 86\% \\
\quad scheduled sampling, $p = 0.35$ & $+0.020$ $\pm$ 0.004 & 11.8 & 0/3 & 0.610 & 87\% \\
\quad first-order Markov chain & $+0.004$ $\pm$ 0.000 & 2.6 & 0/3 & 0.240 & 97\% \\
\quad independent draws & $+0.004$ $\pm$ 0.000 & 2.8 & 0/3 & 0.148 & 97\% \\
\midrule
\multicolumn{6}{@{}l}{\emph{Diagnostic controls: condition on the realised holdout}}\\
\quad arm C: real decisions, lottery re-drawn & $+0.002$ $\pm$ 0.000 & 1.6 & 3/3 & 0.332 & 94\% \\
\quad teacher forcing & $-0.000$ $\pm$ 0.000 & -0.1 & 3/3 & 0.238 & 89\% \\
\bottomrule
\end{tabular}
\end{minipage}\par\bigskip

\noindent\begin{minipage}{\linewidth}\centering\small
\captionof{table}{\textbf{27. Satiation signature, $k = 3+$}: $P(\text{buy on a banner} \mid \text{holding 3+ copies of its featured product})$, pooled over figure A, figure B and weapon. Real-vs-real band $q_{0.95} = 0.001$; real mean 0.146, undefined for 97\% of customers.}
\begin{tabular}{@{}lrrcrr@{}}
\toprule
 & energy & $z$ & inside & gen.\ mean & undefined\\
\midrule
\multicolumn{6}{@{}l}{\emph{Generators a manager could use}}\\
\quad ProductGPT, free running & $+0.025$ $\pm$ 0.005 & 49.2 & 0/3 & 0.683 & 92\% \\
\quad scheduled sampling, $p = 0.15$ & $+0.007$ $\pm$ 0.002 & 14.9 & 0/3 & 0.683 & 96\% \\
\quad scheduled sampling, $p = 0.35$ & $+0.006$ $\pm$ 0.001 & 11.2 & 0/3 & 0.676 & 96\% \\
\quad first-order Markov chain & $+0.000$ $\pm$ 0.000 & 1.0 & 3/3 & 0.241 & 100\% \\
\quad independent draws & $+0.000$ $\pm$ 0.000 & 1.0 & 3/3 & 0.088 & 100\% \\
\midrule
\multicolumn{6}{@{}l}{\emph{Diagnostic controls: condition on the realised holdout}}\\
\quad arm C: real decisions, lottery re-drawn & $+0.000$ $\pm$ 0.000 & 0.6 & 3/3 & 0.271 & 99\% \\
\quad teacher forcing & $-0.000$ $\pm$ 0.000 & -0.3 & 3/3 & 0.179 & 97\% \\
\bottomrule
\end{tabular}
\end{minipage}\par\bigskip


\end{document}
